\section{讨论}\label{sec:discussion}

\subsection{实验实现}

消融结果（第~\ref{sec:ablation}~节）让实现大为简化：不需要逐梳齿的调制器。
一个与微环自由光谱范围锁定的 $200~\mathrm{GHz}$ 输入光梳分成两路，一路保持静态，另一路经过一个承载信号的强度或相位调制器；
两路各过一个只设相位（或幅度和相位）的谱整形器，分别写入式~\eqref{eq:mask} 中的偏置 $b_\mu$ 和权重 $w_\mu$，合束后与泵浦一起送入微环。
商用的硅基液晶谱整形器能独立设定各通道的幅度和相位，最窄通带约 $10~\mathrm{GHz}$、设定步长 $1~\mathrm{GHz}$，远小于梳齿间隔~\cite{coherent_waveshaper}；
C 波段的型号覆盖约 $5~\mathrm{THz}$（1527 到 $1568~\mathrm{nm}$），够 16 根输入梳齿（$3.2~\mathrm{THz}$），32 根以上要更宽的型号；
它的插入损耗（典型约 $5~\mathrm{dB}$，最大 6.5 到 $7~\mathrm{dB}$）要计入输入功率预算。
读出端是下载端口后的解复用器加 31 个光电探测器，结构与频率复用的光学储备池~\cite{butschek2020,lupo2023} 相同。
最难的部分是输入光梳本身：它要与微环的自由光谱范围对齐，并且相位稳定；本文没有模拟它的噪声。

光源的要求来自第~\ref{sec:laser}~节：泵浦激光相对于谐振的本征线宽要在几千赫兹以下。
这排除了普通的 DFB 激光（$100~\mathrm{kHz}$ 到 $1~\mathrm{MHz}$），需要窄线宽激光或自注入锁定~\cite{shen2020}。
慢的频率漂移（远长于一个符号）则要靠锁定环路压到半线宽的约 0.3\% 以内，否则训练好的读出层会随漂移失效（补充材料的闭环实验中，1\% 的慢漂移已让跟踪误差变差 2.4 倍）。

\subsection{功耗与速率}

最好的单环需要约 $180~\mathrm{mW}$ 的总输入光功率（加大下载端口耦合组），其余掩模的最优点在约 20 到 $200~\mathrm{mW}$ 之间；
换成 Si$_3$N$_4$ 微环可以降到约 1/4，但到达探测器的光同样变少，要配更好的探测器或光放大器（第~\ref{sec:platform}~节）。
速率被 Kerr 非线性所需的 $Q$ 限制在约 $1~\mathrm{GBd}$。
以 $1~\mathrm{GBd}$、$180~\mathrm{mW}$ 计，每个符号的光能约 $180~\mathrm{pJ}$，还没有计入激光器的电光转换效率、31 路探测和读出层的计算。
我们没有把它和数字电路做能耗比较，因为那需要一个具体的系统场景，而本文的结论（第~\ref{sec:software}~节）是在精度上微环不占优，能耗上的好处也就只能在信号本来在光域、数字化可以省掉的场景里讨论。

\subsection{局限}

\emph{全部是仿真。} 本文的结论依赖 LLE 模型和噪声模型。我们按规格书和文献选了探测器噪声（每根梳齿 $0.3~\mu\mathrm{W}$，位于 1~GHz 级 InGaAs PIN 接收机 0.13 到 $0.63~\mu\mathrm{W}$ 的中间）、
激光强度噪声和频率噪声，并对后者按真实频谱复核；但微环的热折射噪声、光放大器的自发辐射自拍频、探测器量子效率和输入光梳的噪声都没有建模。
对文献方案的复现也有不确定性：原文没有给出腔长、$Q$ 的分配和采样间隔，我们的对比应理解为“在本文的实现中”。

\emph{读出端。} 读出端的量化只在最好的单环上、用一种标定方式算过（第~\ref{sec:quant}~节）；在光域完成加权的读出没有模拟。

\emph{任务。} 三个基准任务和一个光纤链路任务都是标准任务，掩模的优势集中在需要强非线性、短记忆的预测上；需要长记忆的任务要靠外加延迟线（第~\ref{sec:memory}~节）。

\subsection{试过但不成立的方向}

为了寻找微环相对软件的优势，我们还试了几个方向，结果都是否定或有限的，这里简述，以免后来者重复：
\begin{itemize}
\item \emph{闭环决策}（补充材料）：微环预测 Hénon 映射、再据此控制。微环能在闭环里工作，决策质量由预测噪声决定；
      它只在软件输入噪声取 $3\times10^{-3}$ 时略好于 ESN 和 NG-RC，低于约 $2.4\times10^{-3}$ 时不再领先。
\item \emph{用微环控制另一个 Kerr 微环}：被控对象换成单模 Kerr 腔后，微环控制器的跟踪误差（$0.023$）与 ESN 31（$0.023$）和带前馈的比例微分控制器（$0.022$）相同，
      只比单纯的比例微分控制器（$0.068$）好约 3 倍；稳住不稳定态时，比例微分控制器和 ESN 的成功率（94\% 和 91\%）都高于微环（68\%）。
\item \emph{共参考环}：用一个不加信号的环测出共模的失谐抖动并在读出端减掉。理想参考下也只降 2\%，见第~\ref{sec:reference}~节。
\item \emph{传感}：用同一个光梳估计外加的失谐偏移（以半线宽为单位的均方根误差）。
      窄量程（$\pm0.3$ 个半线宽）内，31 根梳齿的光梳读出只比低功率泵浦加参考的经典读法好约 1.5 倍（$1.0\times10^{-3}$ 对 $1.5\times10^{-3}$）；
      宽量程（$\pm2$ 个半线宽）内光梳远好于只读泵浦（$0.07$ 到 $0.09$ 对 $0.5$ 以上），但几乎没有 Kerr 效应的低功率光梳更好（$0.042$），
      说明好处来自多根梳齿同时探测，而不是 Kerr 非线性。
\end{itemize}
这些结果和第~\ref{sec:software}~节一致：在精度上，Kerr 微梳储备池能做到的是接近一个小型、带量化噪声的软件储备池；
它的价值只能来自物理层（信号本来在光域、延迟、能耗），本文没有证明这一点，我们把它作为后续工作的问题提出。
